\begin{abstract}
Money-laundering detectors are scored per transaction; investigators act on rings. We
give a ring-level evaluation protocol for AMLworld (a ring metric, a matching rule with a
deterministic tie-break, a ring bootstrap that keeps multiplicities, and ring-folded
cross-fitting) and audit the benchmark with it. On HI-Small (test days 8--9, 174 rings, five
seeds), three graph neural detectors spanning test edge F1 0.596 to 0.715 give seeds-only
ring-F1 0.518 to 0.551, and nine post-detector recovery attempts (five read on test, four
stopped at their validation gate) leave ring-F1 flat; most remaining losses are rings with too
few visible members. Below that range the detector still matters: on HI-Medium
(test days 6--14, 2{,}187 rings), a detector trained only on HI-Small, applied without
retraining, gives 0.396 $\pm$ 0.015 (five seeds), against 0.090 for an earlier in-domain detector
that differs in architecture and training data (two seeds). TAML, an account-flagging method
converted to rings, reaches 0.327 [0.254, 0.403] (three seeds), losing mainly through
transactions its graphs omit. We list the evaluation defects we found.
% source: results/notes/final_numbers.md (detector swap rows: edge F1 0.596 / 0.686 / 0.715, B0 0.521 / 0.518 / 0.551;
%   HI-Medium transfer section: 0.396 +- 0.015, 0.090; TAML section: 0.327 [0.254, 0.403]); recovery count: results/notes/recovery_attempts.md;
%   "mostly too few visible members": Q0 bucket a, 65.6 of the 93.8 unmatched rings per seed (results/notes/q_track.md)
\end{abstract}
